\chapter{Estado del arte}
\label{chap:estado_arte}

\ph{Este capítulo sitúa el trabajo en su contexto académico y técnico.
Su función es doble: (1) introducir al lector los conceptos teóricos que necesita
para entender el resto del documento y (2) demostrar que el autor conoce
las soluciones existentes y dónde se sitúa la suya.}

\ph{Estructura recomendada: 3--4 secciones temáticas que cubran los pilares
conceptuales del TFG, más una sección final de \emph{Trabajos relacionados}
donde se comparan soluciones existentes con la propia.}

\ph{Regla de oro de citación: toda definición, dato o afirmación que no sea
tuya lleva su cita con \textbackslash cite junto a la propia afirmación, no
acumulada al final del párrafo. Prefiere la fuente primaria (el artículo o
libro que introdujo el concepto) y combínala con publicaciones recientes
(últimos cinco años) y con tipos de fuente diversos (artículos, congresos,
libros, documentación oficial, software). Patrón de ejemplo: «El concepto~X
fue introducido por Apellido et al.~\cite{apellido2024articulo} y desde
entonces se ha aplicado a \dots». Consulta \texttt{bibliography.bib}: incluye
un modelo de cada tipo de entrada.}

\section{\ph{Concepto / paradigma principal}}
\label{sec:concepto1}

\ph{Primer pilar teórico del trabajo. Por ejemplo: el paradigma, la metodología
o el tipo de sistema sobre el que se construye la solución, apoyado en su
fuente primaria, típicamente un libro de referencia~\cite{apellido2021libro}.}

\ph{Cubre aquí:}
\begin{itemize}
  \item \ph{Definición y origen del concepto.}
  \item \ph{Por qué es relevante para este TFG.}
  \item \ph{Conceptos auxiliares que el lector necesita conocer.}
\end{itemize}

\subsection{\ph{Subconcepto 1.1}}
\label{subsec:concepto1_1}

\ph{Profundiza en una de las ideas del concepto anterior.
Aquí van las definiciones, propiedades o variantes relevantes.}

\subsection{\ph{Subconcepto 1.2}}
\label{subsec:concepto1_2}

\ph{Otro aspecto del mismo concepto. Por ejemplo, su modelo formal,
su notación, su clasificación o sus etapas.}

\subsection{\ph{Herramientas y ecosistema}}
\label{subsec:herramientas1}

\ph{Si procede, repasa las herramientas, librerías o frameworks asociados al concepto
y comenta brevemente cuáles son relevantes para el trabajo y por qué. Cita cada
herramienta con su entrada de software~\cite{autor2024software}.}

\section{\ph{Tecnología / lenguaje / estándar clave}}
\label{sec:concepto2}

\ph{Segundo pilar: una tecnología, lenguaje o estándar concreto sobre el que
se apoya la solución. Sirve para que el lector entienda la sintaxis, la semántica
o los componentes de algo que aparecerá en capítulos posteriores.}

\subsection{\ph{Sintaxis / componentes}}
\label{subsec:sintaxis}

\ph{Describe los elementos básicos.}

\subsection{\ph{Aspecto avanzado / extensión}}
\label{subsec:avanzado}

\ph{Profundiza en un aspecto que será especialmente relevante para el TFG.}

\subsection{\ph{Soporte en herramientas}}
\label{subsec:soporte}

\ph{Si la tecnología tiene un ecosistema de herramientas, comenta las más relevantes
y, si la implementación del TFG decide no usar alguna, explica por qué.}

\section{\ph{Dominio de aplicación}}
\label{sec:dominio}

\ph{Tercer pilar: el dominio físico, funcional o de negocio sobre el que se aplica
la solución. Aporta el conocimiento de dominio mínimo para entender las decisiones
que se tomarán en diseño e implementación, citando los recursos oficiales del
dominio cuando existan~\cite{organizacion2025recurso}.}

\subsection{\ph{Aspecto del dominio 1}}
\label{subsec:dominio_1}

\ph{Por ejemplo: parámetros principales, magnitudes, normativa, actores.}

\subsection{\ph{Aspecto del dominio 2}}
\label{subsec:dominio_2}

\ph{Por ejemplo: limitaciones físicas, criterios de calidad, requisitos legales.}

\section{\ph{Cuarto bloque temático (opcional)}}
\label{sec:concepto4}

\ph{Si el TFG combina más de tres áreas (por ejemplo, una metodología adicional,
un tipo concreto de algoritmo o una tecnología específica), añade aquí un cuarto bloque.
En caso contrario, elimina esta sección.}

\section{Trabajos relacionados}
\label{sec:relacionados}

\ph{Sección clave del estado del arte. Aquí se comparan soluciones, proyectos
o publicaciones que abordan problemas similares al que resuelve el TFG.}

\ph{Empieza con un párrafo que sustente la exhaustividad de la búsqueda:
indica qué bases de datos consultaste (Scopus, Web of Science, Google Scholar,
IEEE Xplore\dots), con qué palabras clave y qué criterios de inclusión y
exclusión aplicaste (p.\,ej.\ publicaciones de los últimos cinco años,
revisadas por pares, con herramienta disponible públicamente). Un párrafo
basta; no hace falta una revisión sistemática formal.}

\ph{Plantilla de redacción:}

\ph{En el ámbito de \dots, destacan los trabajos de \dots, que abordan \dots.
Sin embargo, ninguno cubre \dots, que es precisamente lo que motiva este trabajo.}

\ph{En el ámbito de \dots, los enfoques existentes se centran en \dots.
La diferencia con la propuesta de este TFG es \dots.}

\subsection{Comparativa de trabajos relacionados}
\label{subsec:comparativa_relacionados}

\ph{Resume la comparación en una tabla: una fila por trabajo (citado con
\textbackslash cite) y 3--4 criterios como columnas, derivados del problema
que resuelve tu TFG (no genéricos). La última fila es tu propuesta: así el
lector ve de un vistazo el hueco que cubre el trabajo.}

\begin{table}[H]
  \centering
  \begin{tabular}{P{6cm}ccc}
    \toprule
    \textbf{Trabajo} & \ph{Criterio 1} & \ph{Criterio 2} & \ph{Criterio 3} \\
    \midrule
    \ph{Trabajo A~\cite{apellido2024articulo}} & \ph{Sí} & \ph{Parcial} & \ph{No} \\
    \ph{Trabajo B~\cite{apellido2023congreso}} & \ph{Parcial} & \ph{Sí} & \ph{No} \\
    \ph{Trabajo C~\cite{apellido2022tfm}} & \ph{No} & \ph{Parcial} & \ph{Sí} \\
    \midrule
    \textbf{Este TFG} & \ph{Sí} & \ph{Sí} & \ph{Sí} \\
    \bottomrule
  \end{tabular}
  \caption{Comparativa de trabajos relacionados.}
  \label{tab:comparativa_relacionados}
\end{table}

\ph{Leyenda: Sí = lo cubre completamente; Parcial = lo cubre con limitaciones;
No = no lo cubre. Comenta la tabla de forma crítica, no descriptiva: justifica
cada «Parcial» y cada «No», identifica el hueco que dejan los trabajos
existentes y cierra con una frase que enuncie la contribución diferencial de
este TFG. Sé honesto: si tu propuesta también tiene algún «Parcial»,
reconocerlo refuerza la credibilidad del análisis. Los criterios (columnas)
elegidos deberán poder respaldarse con evidencia propia en el
Capítulo~\ref{chap:resultados}. El Capítulo~\ref{chap:resultados} retomará
esta tabla para contrastarla con los resultados obtenidos.}

% --- Transición al siguiente capítulo ---
\ph{Cierra el capítulo con un párrafo puente de 2--4 líneas: qué queda
establecido aquí y qué pregunta abre el capítulo siguiente. Por ejemplo:
«Identificado el hueco que dejan los trabajos existentes, el
Capítulo~\ref{chap:planificacion} describe cómo se ha organizado el proyecto
para construir la solución que lo cubre: metodología, fases, recursos y
gestión de riesgos».}
